\documentclass[11pt,a4paper]{article}
\usepackage[utf8]{inputenc}
\usepackage[T1]{fontenc}
\usepackage{geometry}
\usepackage{hyperref}
\usepackage{booktabs}
\usepackage{enumitem}
\usepackage{listings}
\usepackage{xcolor}
\geometry{margin=2.5cm}

\lstset{
  basicstyle=\ttfamily\small,
  breaklines=true,
  frame=single,
  backgroundcolor=\color{gray!10},
  numbers=left,
  numberstyle=\tiny
}

\title{Project Report: Open Authorization Protocol}
\author{02244 Logic for Security\\
Fahmida Khan (s250116), Zeinab Ali (s243588), Mithra Arul (s257301)}
\date{March 16, 2026}

\begin{document}
\maketitle
\tableofcontents
\newpage

\section{Open Authorization Protocol -- Zeinab}
\subsection{Description of the final protocol in AnB}
This report presents the final Open Authorization model where user $A$ delegates photo access at server $B$ to service $p$, with trust anchored in identity provider $I$.
The model is analyzed in OFMC and documented together with Week 2--6 refinements.

\subsection{Outline of the design}
\subsubsection{Agents and initial knowledge}
\begin{itemize}[leftmargin=1.2em]
  \item \textbf{$A$ (User):} initiates delegation and proves password knowledge to $I$ using $pw(A,I)$.
  \item \textbf{$B$ (Photo server):} releases photo data only when receiving valid authorization evidence.
  \item \textbf{$p$ (Photo-book service):} requests delegated access and presents IDP grant to $B$.
  \item \textbf{$I$ (Identity provider):} verifies user proof and issues signed grants.
\end{itemize}

For the final hardened model, channel interactions between $A$ and $I$ are represented with pseudonymous secure channels (\texttt{*->*}) and challenge nonce $NI$.

\subsubsection{Security properties}
\begin{itemize}[leftmargin=1.2em]
  \item \texttt{pw(A,I) guessable secret between A,I}
  \item \texttt{fData(Photos) secret between B,p}
  \item \texttt{p authenticates I on fGrant(A,p,B,F,T,Grant)}
  \item \texttt{B authenticates I on fGrant(A,p,B,F,T,Grant)}
  \item \texttt{I authenticates A on NI}
\end{itemize}

\subsection{Protocol design}
Final message flow (used consistently in this report):
\begin{enumerate}[leftmargin=1.4em]
  \item $p \rightarrow A$: request context \texttt{(p,B,F,T,R)}
  \item $A \,*\!\rightarrow\!*\, I$: request context
  \item $I \,*\!\rightarrow\!*\, A$: challenge nonce $NI$
  \item $A \,*\!\rightarrow\!*\, I$: proof \texttt{h(pw(A,I),A,p,B,F,T,R,NI)}
  \item $I \rightarrow p$: signed grant
  \item $p \rightarrow B$: signed grant + scope $F$
  \item $B \rightarrow p$: encrypted photo payload
\end{enumerate}

The password is never used as an encryption key; it appears only inside the authentication hash.

\subsection{Formalization in AnB}
\begin{lstlisting}
Protocol: week6_guessablepassword_hardened

Types:
  Agent A, B, I, p;
  Number F, T, R, NI, Grant, Photos;
  Function pk, pw, h, fGrant, fData;

Knowledge:
  A: A, B, I, p, pk(I), pw(A,I), h, fGrant, fData;
  B: A, B, I, p, pk(B), pk(I), pk(p), inv(pk(B)), fGrant, fData;
  p: A, B, I, p, pk(p), pk(I), inv(pk(p)), fGrant, fData;
  I: A, B, I, p, pk(B), pk(p), pk(I), inv(pk(I)), pw(A,I), h, fGrant, fData;

Actions:
  p -> A: p, B, F, T, R
  A *->* I: A, p, B, F, T, R
  I *->* A: NI
  A *->* I: A, p, B, F, T, R, NI, h(pw(A,I), A, p, B, F, T, R, NI)
  I -> p: {fGrant(A,p,B,F,T,Grant)}inv(pk(I))
  p -> B: {fGrant(A,p,B,F,T,Grant)}inv(pk(I)), F
  B -> p: {fData(Photos)}pk(p)

Goals:
  pw(A,I) guessable secret between A,I;
  fData(Photos) secret between B,p;
  p authenticates I on fGrant(A,p,B,F,T,Grant);
  B authenticates I on fGrant(A,p,B,F,T,Grant);
  I authenticates A on NI;
\end{lstlisting}

\section{Development and Special Tasks (Week 2--4) -- Fahmida}
\subsection{Outline of the design}
Weeks 2--4 build the protocol from a minimal baseline toward structured IDP-based authorization.
The focus is on (i) baseline attacker reasoning, (ii) lazy intruder executability, and (iii) type-flaw resistance with explicit key retrieval.

\subsection{Protocol development and OFMC findings}
\subsubsection{Week 2: Dolev--Yao and static analysis}
Week 2 uses a simplified baseline and passive static reasoning.
The intruder observes visible fields but cannot derive honest private keys or decrypt ciphertext for uncompromised recipients.
This baseline is intentionally weak and serves as evidence for why stronger authorization structure is needed.

\subsubsection{Week 3: Lazy intruder}
Week 3 introduces identity-provider-issued authorization and the special-task executability check.
We verified that role $P$ is executable when $P$ is instantiated by the intruder:
the intruder can construct each required outgoing message from previously available knowledge.
Operationally, the intruder can receive a token from $A$, forward it to $B$, and continue protocol interaction.
This is expected for a delegated service role.
The intended security condition is not ``intruder cannot send,'' but rather ``only correctly issued IDP-bound authorization should be accepted by $B$.''
Compared with the Week 2 baseline, Week 3 improves the protocol by introducing the identity provider $I$ as the authorization issuer instead of relying on direct delegation trust.
User $A$ now provides explicit password-based authentication evidence through $h(pw(A,I),NA)$, so the password is used for proof of knowledge and not as an encryption key.
The signed grant binds the delegation context (participants and request parameters), which gives $B$ a stronger basis for accepting access requests.
We also completed the special-task lazy-intruder check and showed that role $P$ remains executable even when instantiated by the intruder, which confirms model executability.
However, this executability result does not imply full security, and remaining attack findings motivated the Week 4 refinements (typed formats and key lookup).

\subsubsection{Week 4: Type-flaw resistance and GetPublicKey}
In Week 4 we introduced explicit format constructors (e.g., request/token/data wrappers) to reduce message-type confusion.
The core typing claim is: for any two non-variable message patterns, if they unify, they belong to the same intended kind/type.
We also added a separate GetPublicKey protocol so that $A$, initially trusting only $pk(IDP)$, can obtain $pk(X)$ from the identity provider.
This removes the unrealistic assumption that $A$ initially knows every party public key.

\subsection{AnB formalization (representative snippets)}
\subsubsection{Design decisions and development}
The early versions used stronger initial assumptions to get an executable baseline quickly.
During development, these assumptions were reduced:
\begin{itemize}[leftmargin=1.2em]
  \item direct trust assumptions were replaced by IDP-issued authorization,
  \item typed formats were added to improve message discipline,
  \item key lookup was split into a dedicated protocol before main delegation flow when needed.
\end{itemize}
This staged approach made attack traces and modeling limitations visible at each step and allowed targeted refinement.

\subsubsection{Week 2 snippet (baseline)}
\begin{lstlisting}
Protocol: OpenAuth_W2_Simple
...
Actions:
  A -> P: {A,B,P,F,T}inv(pk(A))
  P -> B: {A,B,P,F,T}inv(pk(A)),F
  B -> P: {Photos}pk(P)
Goals:
  Photos secret between B,P;
  B authenticates A on P,F,T;
\end{lstlisting}

\subsubsection{Week 3 snippet (IDP-based delegation)}
\begin{lstlisting}
Protocol: OpenAuth_W3_DelegationViaIDP
...
Actions:
  P -> A: P,B,F,T,NA
  A -> I: A,P,B,F,T,NA,h(pw(A,I),NA)
  I -> A: {A,P,B,F,T,NA,Grant}inv(pk(I))
  A -> P: {A,P,B,F,T,NA,Grant}inv(pk(I))
  P -> B: {A,P,B,F,T,NA,Grant}inv(pk(I)),F
Goals:
  Photos secret between B,P;
\end{lstlisting}

\subsubsection{Week 4 snippets (typed formats and key lookup)}
\begin{lstlisting}
Protocol: OpenAuth_W4_Formats_Delegation
...
Function pk,pw,h,fReq,fAuthProof,fGrant,fUse,fData;
Actions:
  p -> A: fReq(p,B)
  A -> I: fAuthProof(A,p,B,h(pw(A,I),R))
  I -> A: {fGrant(fAuthProof(A,p,B,h(pw(A,I),R)),Grant)}inv(pk(I))
  p -> B: {fGrant(...)}inv(pk(I)),fUse(fReq(p,B))
Goals:
  fData(Photos) secret between B,p;
\end{lstlisting}

\begin{lstlisting}
Protocol: OpenAuth_W4_KeyLookup
...
Actions:
  A -> I: fKeyReq(A, X)
  I -> A: {fKeyResp(X, pk(X))}inv(pk(I))

Goals:
  A authenticates I on fKeyResp(X, pk(X));
\end{lstlisting}

\section{TLS Channels, Guessable Password, and Final Protocol -- Mithra}
\subsection{Week 5: TLS channel abstraction}
Week 5 models pseudonymous channels for transport abstraction.
OFMC still reports attacks on some variants, showing that channel abstraction alone does not solve logic-level authorization flaws.

\subsection{Week 6: guessable-password analysis}
Week 6 focuses on low-entropy credential robustness.
Results are interpreted cautiously: for the dedicated guessable-password goal and tested bounds, no guessing attack is reported in the focused model.
This should not be generalized to unconditional security of all goals and all session bounds.

\subsection{Final assumptions and remaining limitations}
\begin{itemize}[leftmargin=1.2em]
  \item Trusted IDP assumption is central.
  \item Resource and photo payload are abstract symbolic terms.
  \item Bounded-session verification does not imply full unbounded guarantees.
  \item Different variants may satisfy different subsets of goals.
\end{itemize}

\subsection{Summary of special tasks}
\begin{center}
\begin{tabular}{ll}
\toprule
\textbf{Week} & \textbf{Outcome} \\
\midrule
2 & Static Dolev--Yao analysis completed on baseline model. \\
3 & Lazy intruder executability of role $p$ demonstrated. \\
4 & Type-flaw-resistant format strategy and key-lookup protocol documented. \\
5 & Pseudonymous channel abstraction implemented and analyzed. \\
6 & Guessable-password checks completed with focused bounded no-attack result. \\
\bottomrule
\end{tabular}
\end{center}

\subsection{AnB files delivered}
\begin{itemize}[leftmargin=1.2em]
  \item \texttt{Week2\_Problem2\_First\_AnB\_Version.anb}
  \item \texttt{Week3\_Problem1\_AnB\_With\_IdentityProvider.anb}
  \item \texttt{Week3\_Problem3\_Refined\_HonestProvider.anb}
  \item \texttt{Week4\_Problem1\_MainProtocol\_Formats.anb}
  \item \texttt{Week4\_Problem2\_AnB\_KeyLookup\_Protocol.anb}
  \item \texttt{Week5\_Problem1\_PseudonymousChannels.anb}
  \item \texttt{Week6\_Problem1\_GuessablePassword\_Check.anb}
  \item \texttt{Week6\_Problem2\_GuessablePassword\_Hardened.anb}
  \item \texttt{Week6\_Problem3\_GuessablePassword\_OnlyCheck.anb}
\end{itemize}

\appendix
\section{Week 3 Full AnB Files}

\subsection{Week3\_Problem1\_AnB\_With\_IdentityProvider.anb}
\begin{lstlisting}[caption={Week 3 -- Initial IDP integration model},label={lst:week3-idp-full}]
Protocol: OpenAuth_W3_DelegationViaIDP

Types:
  Agent A,B,P,I;
  Number F,T,NA,Grant,Photos;
  Function pk,pw,h;

Knowledge:
  A: A,B,P,I,pk(A),pk(B),pk(P),pk(I),inv(pk(A)),pw(A,I),h;
  B: A,B,P,I,pk(A),pk(B),pk(P),pk(I),inv(pk(B));
  P: A,B,P,I,pk(A),pk(B),pk(P),pk(I),inv(pk(P));
  I: A,B,P,I,pk(A),pk(B),pk(P),pk(I),inv(pk(I)),pw(A,I),h;

Actions:
  P -> A: P,B,F,T,NA
  A -> I: A,P,B,F,T,NA,h(pw(A,I),NA)
  I -> A: {A,P,B,F,T,NA,Grant}inv(pk(I))
  A -> P: {A,P,B,F,T,NA,Grant}inv(pk(I))
  P -> B: {A,P,B,F,T,NA,Grant}inv(pk(I)),F
  B -> P: {Photos}pk(P)

Goals:
  Photos secret between B,P;
  B authenticates I on A,P,F,T,Grant;
\end{lstlisting}

\subsection{Week3\_Problem3\_Refined\_HonestProvider.anb}
\begin{lstlisting}[caption={Week 3 -- Refined honest-provider model},label={lst:week3-refined-full}]
Protocol: OpenAuth_W3_Refined_HonestProvider

Types:
  Agent A,B,I,p;
  Number F,T,NA,Grant,Photos;
  Function pk,pw,h;

Knowledge:
  A: A,B,p,I,pk(A),pk(B),pk(p),pk(I),inv(pk(A)),pw(A,I),h;
  B: A,B,p,I,pk(A),pk(B),pk(p),pk(I),inv(pk(B));
  p: A,B,p,I,pk(A),pk(B),pk(p),pk(I),inv(pk(p));
  I: A,B,p,I,pk(A),pk(B),pk(p),pk(I),inv(pk(I)),pw(A,I),h;

Actions:
  p -> A: p,B,F,T,NA
  A -> I: A,p,B,F,T,NA,h(pw(A,I),NA)
  I -> A: {A,p,B,F,T,NA,Grant}inv(pk(I))
  A -> p: {A,p,B,F,T,NA,Grant}inv(pk(I))
  p -> B: {A,p,B,F,T,NA,Grant}inv(pk(I)),F
  B -> p: {Photos}pk(p)

Goals:
  Photos secret between B,p;
  B authenticates I on A,p,F,T,Grant;
  p authenticates B on Photos;
\end{lstlisting}

\section{Week 4 Full AnB Files}

\subsection{Week4\_Problem1\_MainProtocol\_Formats.anb}
\begin{lstlisting}[caption={Week 4 -- Main protocol with typed formats},label={lst:week4-formats-full}]
Protocol: OpenAuth_W4_Formats_Delegation

Types:
  Agent A,B,I,p;
  Number R,Grant,Photos;
  Function pk,pw,h,fReq,fAuthProof,fGrant,fUse,fData;

Knowledge:
  A: A,B,p,I,pk(A),pk(B),pk(p),pk(I),inv(pk(A)),pw(A,I),h,fReq,fAuthProof,fGrant,fUse,fData;
  B: A,B,p,I,pk(A),pk(B),pk(p),pk(I),inv(pk(B)),fReq,fAuthProof,fGrant,fUse,fData;
  p: A,B,p,I,pk(A),pk(B),pk(p),pk(I),inv(pk(p)),fReq,fAuthProof,fGrant,fUse,fData;
  I: A,B,p,I,pk(A),pk(B),pk(p),pk(I),inv(pk(I)),pw(A,I),h,fReq,fAuthProof,fGrant,fUse,fData;

Actions:
  p -> A: fReq(p,B)
  A -> I: fAuthProof(A,p,B,h(pw(A,I),R))
  I -> A: {fGrant(fAuthProof(A,p,B,h(pw(A,I),R)),Grant)}inv(pk(I))
  A -> p: {fGrant(fAuthProof(A,p,B,h(pw(A,I),R)),Grant)}inv(pk(I))
  p -> B: {fGrant(fAuthProof(A,p,B,h(pw(A,I),R)),Grant)}inv(pk(I)),fUse(fReq(p,B))
  B -> p: {fData(Photos)}pk(p)

Goals:
  fData(Photos) secret between B,p;
  B authenticates I on {fGrant(fAuthProof(A,p,B,h(pw(A,I),R)),Grant)}inv(pk(I));
\end{lstlisting}

\subsection{Week4\_Problem2\_AnB\_KeyLookup\_Protocol.anb}
\begin{lstlisting}[caption={Week 4 -- GetPublicKey protocol},label={lst:week4-keylookup-full}]
Protocol: OpenAuth_W4_KeyLookup

Types:
  Agent A,I,X;
  Function pk,fKeyReq,fKeyResp;

Knowledge:
  A: A,I,X,pk(A),pk(I),inv(pk(A)),fKeyReq,fKeyResp;
  I: A,I,X,pk(A),pk(I),pk(X),inv(pk(I)),fKeyReq,fKeyResp;

Actions:
  A -> I: fKeyReq(A,X)
  I -> A: {fKeyResp(X,pk(X))}inv(pk(I))

Goals:
  A authenticates I on fKeyResp(X,pk(X));
\end{lstlisting}

\section*{References}
\begin{itemize}[leftmargin=1.2em]
  \item Sebastian M\"odersheim, \emph{Protocol Security Verification Tutorial}, DTU Compute.
  \item OFMC and AnB documentation.
  \item 02244 course slides and assignment handout.
\end{itemize}

\end{document}
